\section{Related Work}\label{s:related}
Simulation-based testing is standard practice in the aerospace sector, where real flight tests are expensive, risky, and difficult to repeat under controlled conditions~\cite{bittar-guidance-sil-uav-xplane-simulink, santos-hil-uav-autopilot}. The dominant approach couples flight software with a physics-based simulator, exchanging states throughout the simulation. This pattern is particularly well-established in the Unmanned Aerial Vehicle (UAV) sector. 

Santos et al.~\cite{santos-hil-uav-autopilot} built an HIL environment that couples a PC flight simulator, X-Plane, with a dedicated microcontroller running a full lateral and longitudinal autopilot. Since X-Plane operates natively in real time, the two systems run independently, communicating over Ethernet using UDP at 20 Hz. Santos et al. validated the HIL framework by testing longitudinal and lateral control loops on a PT-60 radio-controlled aircraft model in X-Plane. Both control loops within the framework achieved the desired altitude and roll behaviour throughout the respective operations. The findings of their study confirm that the HIL environment is sufficient for developing and tuning control parameters without physically flying the aircraft and risking the airframe during testing.


% For the longitudinal axis, the autopilot was commanded to climb from an initial altitude to 3000ft, then to 5000ft. In both cases, it reached the desired altitude and demonstrated a stable level flight. For the lateral axis, the autopilot was commanded to change heading from 120\degree~to 200\degree~, and achieved stable roll behaviour throughout the operation. 


Bittar et al.~\cite{bittar-guidance-sil-uav-xplane-simulink} follow the same pattern, but transition to a SIL context, coupling X-Plane as the flight dynamics simulator with Simulink executing the control algorithms, communicating over UDP at 20Hz, on the same machine. These studies show that this independent-process pattern holds regardless of whether flight software runs on dedicated hardware or on a host computer.

% \begin{figure}[h]
%     \centering
%     \includegraphics[width=0.75\linewidth]{figures/lit_review/matlab-xplane-bittar.png}
%     \caption{SIL architecture showcasing Simulink executing the guidance and control algorithms while X-Plane handles the vehicle dynamics, with state exchanged over UDP. Reproduced from Bittar et al.~\cite{bittar-guidance-sil-uav-xplane-simulink}.}
%     \label{fig:matlab-xplane-bittar}
% \end{figure}

Coopmans et al.~\cite{coopmans-sil-hil-verification-uav} further validate that simulated UAV trajectories using SIL and HIL setups closely match real flight data. The authors demonstrate this using JSBSim as the flight simulator and the open-source Paparazzi autopilot as the flight software. In SIL mode, the Paparazzi control code is compiled and executed on a laptop against the JSBSim model, without additional hardware required. In HIL mode, the same code runs unchanged on the actual autopilot hardware, with JSBSim providing simulated sensor input through a serial driver in place of physical sensors. In their tests, the flight plan remained the same across all three modes: SIL, HIL and real flight, enabling a direct comparison. All three modes produced similar trajectories, remaining within 10\% error despite 6~$m/s$ wind conditions, leading the authors to conclude that both SIL and HIL are sufficiently accurate for rapid development and tuning of control parameters.

% \begin{figure}[h]
%     \centering
%     \includegraphics[width=0.75\linewidth]{figures/lit_review/figure-8-coopman.png}
%     \caption{UAV Trajectories of a figure-eight manoeuvre across SIL, HIL and real flight, demonstrating close agreement between all modes. Reproduced from Coopmans et al.~\cite{coopmans-sil-hil-verification-uav}.}
%     \label{fig:figure-8-coopman}
% \end{figure}

This pattern extends beyond UAVs to small spacecraft, a domain closer to the context of high-powered rockets addressed in this study. Kiesbye et al.~\cite{kiesbye2019-moveII-cubsat} developed a combined SIL and HIL environment for the MOVE-II CubeSat at the Technical University of Munich. The entire simulation, including the space environment, orbital dynamics, sensor models and actuator models, runs in Matlab/Simulink. In SIL mode, the attitude control algorithms are implemented in C++, embedded into the Simulink model as Level-2 S-functions, and executed within the simulation environment. In HIL mode, the flight code runs on the satellite hardware, with Simulink transmitting simulated sensor data via interface simulators over UDP, enabling communication with the hardware via its Serial Peripheral Interface (SPI) bus, and reading actuator commands back to close the loop. The authors demonstrate that this simulation environment supports the development and verification of MOVE-II's flight software throughout its mission, and continues to be used after launch to analyse in-flight issues and verify software updates before uplink. Notably, the SIL approach requires embedding flight logic as S-function blocks within Simulink, a constraint that does not extend to flight software frameworks such as Paparazzi and F Prime, which run as independent processes with their own threading models and inter-process communication.



The only published work connecting RocketPy to embedded hardware is Netzel et al. \cite{netzel-rocketpy-hil}, who validated their onboard flight software apogee detection algorithm by replaying pre-computed RocketPy trajectories as sensor data over SPI. While this confirms that the hardware correctly processes sensor readings, the approach is open-loop, as the simulation is pre-computed and not influenced by the outputs of the flight software. Control algorithms and recovery timing under dynamic flight conditions, therefore, cannot be validated using this method. RITL addresses this limitation directly by coupling RocketPy and F Prime in a live closed feedback loop, where actuation outputs influence the simulation state at each step.

These works establish simulation-based verification as a viable and well-validated approach for aerospace software, spanning UAV autopilots, CubeSat attitude control, and rocket apogee detection. However, none provides a closed-loop coupling framework between a high-powered rocketry trajectory simulator and a flight software framework. Existing HPR tools remain open-loop, and existing SIL/HIL frameworks are tightly coupled to specific simulation environments, such as Matlab/Simulink. RITL addresses this gap by coupling RocketPy with F Prime in a closed feedback loop, supporting SIL and HIL execution and enabling fault injection within the simulation loop.



